\documentclass[11pt]{article}
\usepackage[margin=1in]{geometry}
\usepackage{amsmath,amssymb}
\usepackage{algorithm}
\usepackage{algpseudocode}
\usepackage{fancyhdr}
\usepackage{enumerate}
\usepackage[colorlinks=true,linkcolor=blue,urlcolor=blue]{hyperref}

\newcommand{\hwnum}{7}

% page header, same as the posted solutions
\pagestyle{fancy}
\fancyhf{}
\lhead{EECS 376: Foundations of Computer Science\\University of Michigan, Fall 2026}
\rhead{\vspace{\baselineskip}Solutions to Homework \hwnum}
\cfoot{\thepage}
\setlength{\headheight}{28pt}

% pseudocode labels
\renewcommand{\algorithmicrequire}{\textbf{Input:}}
\renewcommand{\algorithmicensure}{\textbf{Output:}}

% notation
\newcommand{\set}[1]{\left\{#1\right\}}
\newcommand{\abs}[1]{\left|#1\right|}
\newcommand{\inner}[1]{\langle #1 \rangle}
\newcommand{\algo}[1]{\textsc{#1}}
\newcommand{\lT}{\leq_T}
\newcommand{\Lh}{L_{\text{HALT}}}
\newcommand{\Lacc}{L_{\text{ACC}}}
\newcommand{\Lb}{L_{\text{BARBER}}}
\newcommand{\Leacc}{L_{\varepsilon\text{-ACC}}}
\newcommand{\Lrec}{L_{\textsc{Rec}}}
\newcommand{\Lbh}{L_{\text{BOUNDED-HALT}}}
\newcommand{\Llh}{L_{\text{LINEAR-HALT}}}
\newcommand{\Lloops}{L_{\text{LOOPS}}}
\newcommand{\LSp}{L_{S+}}
\newcommand{\Llag}{L_{\text{LAG}}}
\newcommand{\Lsmh}{L_{\text{SmallHalt}}}
\newcommand{\Ms}{M_{\textsc{Search}}}
\newcommand{\Mb}{M_{\text{BARBER}}}

% one part per page for Gradescope
\newcommand{\qpage}[1]{\clearpage\noindent\textbf{\large #1}\par\medskip}

\begin{document}

\noindent\textbf{\large Q0 Before you start; before you submit}\par\medskip

I read Handout 5 on Turing reductions before starting, and the reductions below follow it: name the oracle, give pseudocode for the machine that uses it, show that the machine always halts, and prove correctness in both directions.
I did not use Rice's Theorem anywhere.
I did not collaborate with anyone on this homework.

\qpage{Q1 Self assessment}

The part of HW6 with the most room for improvement is Question 3(c), about replacing the oracle $H$ by the machine $M_{SH}$ that decides the finite language $\Lsmh$.
Here is what I submitted:

\begin{quote}\small
Replacing $H$ by $M_{SH}$ gives an ordinary Turing machine $M_G'$ (no oracle) that makes one call $M_{SH}(\inner{\Ms}, \varepsilon)$ and returns the opposite answer.
Whether this ``works'' depends on the length of that query.

\textbf{If $(\inner{\Ms}, \varepsilon)$ has $1000$ or more characters}, it is not in $\Lsmh$ by definition, whatever $\Ms$ does.
So $M_{SH}$ rejects it and $M_G'$ always accepts, regardless of whether the conjecture is true.
In this case the solution does not work: $M_G'$'s answer is not based on the behavior of $\Ms$ at all.
(The correctness chain from part (b) breaks at the step ``$(\inner{\Ms}, \varepsilon) \in \Lh \Longrightarrow$ oracle accepts,'' which is no longer justified.)

\textbf{If $(\inner{\Ms}, \varepsilon)$ has fewer than $1000$ characters}, then $M_{SH}$ agrees with $H$ on this particular query, so the proof from part (b) goes through word for word, and $M_G'$ is a Turing machine with no oracle that accepts on every input if the conjecture is true and rejects on every input otherwise.

This sounds exciting, but it does not help us learn the answer.
The reason is that the proof that $\Lsmh$ is decidable is non-constructive.
$\Lsmh$ is finite, so it is decided by a machine that hard-codes the finite list of its members (as in HW5, Q6(a)), but we are not told which strings are on that list.
Writing down $M_{SH}$ would require knowing, for each of the finitely many pairs of length $< 1000$, whether $M$ halts on $x$, and for the pair $(\inner{\Ms}, \varepsilon)$ that is exactly the question of whether the Goldbach Conjecture is true.
So $M_{SH}$ exists, but nobody can actually write it down and run it without already having resolved the conjecture.

In fact, the existence of a Turing machine that answers the Goldbach question was never in doubt: one of the two trivial machines ``always accept'' and ``always reject'' is correct, we just do not know which.
Part (b) is interesting only because the oracle $H$ is a single machine that would answer this and countless other open questions at once, and part (c) shows that the finite restriction $\Lsmh$ does not give us anything like that.
\end{quote}

What was deficient: I never actually answered the question.
The problem asks whether the solution from part (b) still works, and the posted solution answers in two sentences: the only query is $(\inner{\Ms}, \varepsilon)$, it has fewer than $1000$ characters, so $M_{SH}$ returns the same answer as $H$ on it and the whole correctness analysis from part (b) goes through.
I instead split into two cases by the length of the query and treated both as equally plausible, so the real answer is buried in the middle of a case analysis and the rubric item ``state that the query has fewer than $1000$ characters'' is not clearly there.
I also never said why the query is short.
Then I spent the second half on the remark about why this does not settle the conjecture, which is interesting but is not what was asked.

How to fix it: start with the answer.
Yes, it still works.
$\inner{\Ms}$ is the encoding of the short program from part (a) together with a short \algo{isPrime} subroutine, so even with a generous encoding it is a few hundred characters, and $\varepsilon$ adds nothing, so the pair is well under $1000$ characters.
Hence $(\inner{\Ms}, \varepsilon) \in \Lsmh$ if and only if $\Ms(\varepsilon)$ halts, which means $M_{SH}$ and $H$ agree on the only query $M_G$ makes.
The four-line implication chains from part (b) are then unchanged with $M_{SH}$ in place of $H$, and $M_G'$ is an ordinary Turing machine that accepts if and only if the Goldbach Conjecture is true.
The remark that we do not know what $M_{SH}$ looks like can stay, but as a two-sentence note at the end, not as half of the answer.

\qpage{Q2.1 Warm-up exercises (a)}

\textbf{No.} $L_1 \cap L_2$ is not necessarily decidable.

Counterexample: $L_1 = \Lh$ and $L_2 = \Sigma^*$.

\textbf{$L_1 = \Lh$ is recognizable.}
This was shown in lecture.
The recognizer is the universal machine: on input $(\inner{M}, x)$, run $M$ on $x$, and accept as soon as $M$ halts.
If $M(x)$ halts, the simulation finishes and the machine accepts; if $M(x)$ loops, the simulation never finishes, so the machine does not accept.
So it accepts exactly the strings in $\Lh$.

\textbf{$L_2 = \Sigma^*$ is decidable.}
The machine that accepts every input immediately halts on every input and accepts exactly the strings in $\Sigma^*$.

\textbf{$L_1 \cap L_2$ is undecidable.}
Since $\Lh \subseteq \Sigma^*$, we have $L_1 \cap L_2 = \Lh \cap \Sigma^* = \Lh$, and $\Lh$ is undecidable (proved in lecture).

So $L_1$ is recognizable, $L_2$ is decidable, and $L_1 \cap L_2$ is undecidable.
This disproves the claim.

The reason this can happen: intersecting with a decidable language can only remove strings, and intersecting with $\Sigma^*$ removes nothing, so the undecidable part of $L_1$ is still there.
(Part (b) shows that recognizability, on the other hand, does survive.)

\qpage{Q2.2 Warm-up exercises (b)}

\textbf{Yes.} $L_1 \cap L_2$ is always recognizable.

Since $L_1$ is recognizable, there is a Turing machine $R_1$ that accepts exactly the strings in $L_1$ (it may reject or loop on the others).
Since $L_2$ is decidable, there is a Turing machine $D_2$ that halts on every input and accepts exactly the strings in $L_2$.
Consider the following machine.

\begin{algorithmic}[1]
  \Require{A string $x$}
  \Ensure{accept if $x \in L_1 \cap L_2$; otherwise reject or loop}
  \Function{$R$}{$x$}
    \If{$D_2(x)$ rejects} \Return reject
    \EndIf
    \State Run $R_1(x)$ \Comment{this may loop, which is allowed for a recognizer}
    \If{$R_1(x)$ accepts} \Return accept
    \Else{} \Return reject
    \EndIf
  \EndFunction
\end{algorithmic}

A recognizer does not have to halt on every input, so the only thing to prove is that $R$ accepts $x$ if and only if $x \in L_1 \cap L_2$.

\textbf{If $x \in L_1 \cap L_2$, then $R$ accepts $x$.}
Since $x \in L_2$ and $D_2$ decides $L_2$, the first call $D_2(x)$ halts and accepts, so the first \textbf{if} does not fire.
Since $x \in L_1$ and $R_1$ recognizes $L_1$, $R_1(x)$ halts and accepts.
So $R$ reaches the last \textbf{if} and accepts.

\textbf{If $x \notin L_1 \cap L_2$, then $R$ does not accept $x$.}
There are two cases.
\begin{itemize}
  \item $x \notin L_2$. Then $D_2(x)$ halts and rejects, so $R$ rejects on the first \textbf{if}.
  \item $x \in L_2$ but $x \notin L_1$. Then $D_2(x)$ accepts and $R$ goes on to run $R_1(x)$.
  Since $x \notin L_1$, $R_1(x)$ does not accept: it either rejects, in which case $R$ rejects, or it loops, in which case $R$ loops.
  Either way $R$ does not accept.
\end{itemize}
So $R$ accepts exactly the strings in $L_1 \cap L_2$, i.e., $R$ recognizes $L_1 \cap L_2$.

(Calling $D_2$ first is not essential, since a recognizer may loop anyway, but it is the natural order: the decidable test is the one that is guaranteed to finish.)

\qpage{Q2.3 Warm-up exercises (c)}

The key observation is that \emph{every} Turing machine recognizes some language.
By definition, $M$ recognizes $L$ if $M$ accepts exactly the strings in $L$.
Given any Turing machine $M$, let $L(M) = \set{x : M(x) \text{ accepts}}$.
This is a perfectly good language, and $M$ accepts exactly the strings in it, so $M$ recognizes $L(M)$.
It does not matter whether $M$ loops on some inputs, since a recognizer is allowed to loop on strings outside its language.
So
\[
  \Lrec = \set{\inner{M} : M \text{ is a Turing machine}},
\]
the set of all strings that are valid encodings of Turing machines.

As in lecture, a valid encoding $\inner{M}$ is a string that lists the components of $M$: the finite set of states, the input and tape alphabets, the transition function as a finite table, and the start, accept, and reject states.
Checking whether a given string has this form is a finite syntactic check, with no need to run $M$.

\begin{algorithmic}[1]
  \Require{A string $w$}
  \Ensure{accept if $w \in \Lrec$, reject otherwise}
  \Function{$D_{\textsc{Rec}}$}{$w$}
    \State Try to parse $w$ as an encoding $(Q, \Sigma, \Gamma, \delta, q_0, q_{\text{acc}}, q_{\text{rej}})$ of a Turing machine
    \State Check that $Q, \Sigma, \Gamma$ are finite lists, that $\delta$ assigns to each pair in $Q \times \Gamma$ exactly one triple in $Q \times \Gamma \times \set{L, R}$, and that $q_0, q_{\text{acc}}, q_{\text{rej}} \in Q$
    \If{every check passes} \Return accept
    \Else{} \Return reject
    \EndIf
  \EndFunction
\end{algorithmic}

\textbf{$D_{\textsc{Rec}}$ always halts.}
Its input $w$ is a finite string, and every step is a finite check on that string: parsing it, counting entries, and comparing the entries of the transition table against the finite lists of states and symbols.
$D_{\textsc{Rec}}$ never simulates the machine described by $w$.

\textbf{Correctness.}
\begin{itemize}
  \item If $w \in \Lrec$, then $w = \inner{M}$ for some Turing machine $M$ (which recognizes $L(M)$), so $w$ is a valid encoding, every check passes, and $D_{\textsc{Rec}}$ accepts.
  \item If $w \notin \Lrec$, then $w$ is not the encoding of any Turing machine: if it were, that machine would recognize its own language $L(M)$ and $w$ would be in $\Lrec$.
  So some check fails and $D_{\textsc{Rec}}$ rejects.
\end{itemize}
Hence $D_{\textsc{Rec}}$ decides $\Lrec$, and $\Lrec$ is decidable.

(Under the alternative convention where every string is treated as the encoding of some Turing machine, say a fixed default machine for malformed strings, $\Lrec$ is just $\Sigma^*$ and the decider is ``always accept.'')

\qpage{Q3.1 Deciding Undecidability (a)}

$\Leacc = \set{\inner{M} : M \text{ accepts } \varepsilon}$ is \textbf{undecidable}.
I show $\Lacc \lT \Leacc$; since $\Lacc$ is undecidable (lecture), so is $\Leacc$.

Let $O$ be an oracle that decides $\Leacc$: $O$ halts on every input, and $O(\inner{N})$ accepts if and only if $N(\varepsilon)$ accepts.
The idea is to turn the question ``does $M$ accept $x$?'' into a question about a machine $M'$ that behaves on every input, in particular on $\varepsilon$, the way $M$ behaves on $x$.

\begin{algorithmic}[1]
  \Require{$(\inner{M}, x)$, where $M$ is a Turing machine and $x$ a string}
  \Ensure{accept if $M(x)$ accepts, reject otherwise}
  \Function{$M_{\text{ACC}}$}{$\inner{M}, x$}
    \State Construct the code $\inner{M'}$ of the following machine $M'$:
    \Statex \hspace{2em} \textbf{function} $M'(w)$: \Comment{$w$ is ignored}
    \Statex \hspace{4em} Run $M(x)$ \Comment{$\inner{M}$ and $x$ are hard-coded into $M'$}
    \Statex \hspace{4em} \textbf{if} $M(x)$ accepts \textbf{then return} accept \textbf{else return} reject
    \State \Return $O(\inner{M'})$
  \EndFunction
\end{algorithmic}

\textbf{$M_{\text{ACC}}$ always halts.}
Constructing $\inner{M'}$ is string manipulation: $M_{\text{ACC}}$ writes a fixed template of code with the strings $\inner{M}$ and $x$ pasted in.
It does not run $M$ or $M'$.
Then it makes one call to $O$, which halts because $O$ decides $\Leacc$, and returns that answer.

\textbf{Key property of $M'$.}
For every input $w$, $M'(w)$ ignores $w$ and runs $M(x)$.
If $M(x)$ accepts, $M'$ accepts.
If $M(x)$ rejects, $M'$ rejects.
If $M(x)$ loops, $M'$ never gets past the simulation, so $M'$ loops.
In particular, for $w = \varepsilon$:
\[
  M'(\varepsilon) \text{ accepts} \iff M(x) \text{ accepts}.
\]

\textbf{Correctness.}
By the definition of $\Lacc$, the key property, the definition of $\Leacc$, the fact that $O$ decides $\Leacc$, and the code of $M_{\text{ACC}}$,
\begin{align*}
  (\inner{M}, x) \in \Lacc &\Longrightarrow M(x) \text{ accepts} \\
  &\Longrightarrow M'(\varepsilon) \text{ accepts} \\
  &\Longrightarrow \inner{M'} \in \Leacc \\
  &\Longrightarrow O(\inner{M'}) \text{ accepts} \\
  &\Longrightarrow M_{\text{ACC}} \text{ accepts}, \\[0.5em]
  (\inner{M}, x) \notin \Lacc &\Longrightarrow M(x) \text{ rejects or loops} \\
  &\Longrightarrow M'(\varepsilon) \text{ rejects or loops, so does not accept} \\
  &\Longrightarrow \inner{M'} \notin \Leacc \\
  &\Longrightarrow O(\inner{M'}) \text{ rejects} \\
  &\Longrightarrow M_{\text{ACC}} \text{ rejects}.
\end{align*}
So $M_{\text{ACC}}$ decides $\Lacc$ using an oracle for $\Leacc$, that is, $\Lacc \lT \Leacc$.
If $\Leacc$ were decidable, we could replace $O$ by its decider and get a decider for $\Lacc$, which does not exist.
Therefore $\Leacc$ is undecidable.

\qpage{Q3.2 Deciding Undecidability (b)}

$L = \set{(\inner{M_1}, \inner{M_2}) : \varepsilon \in L(M_1) \cup L(M_2)}$ is \textbf{undecidable}.
Here $L(M_i)$ is the set of strings $M_i$ accepts, so $(\inner{M_1}, \inner{M_2}) \in L$ exactly when $M_1(\varepsilon)$ accepts or $M_2(\varepsilon)$ accepts (or both).
I show $\Lacc \lT L$; since $\Lacc$ is undecidable, so is $L$.

Let $O$ be an oracle that decides $L$: $O$ halts on every input, and $O(\inner{M_1}, \inner{M_2})$ accepts if and only if $\varepsilon \in L(M_1) \cup L(M_2)$.
The trick is that nothing stops us from handing the oracle the same machine twice: $L(M') \cup L(M') = L(M')$, so the oracle then just tells us whether $M'$ accepts $\varepsilon$.

\begin{algorithmic}[1]
  \Require{$(\inner{M}, x)$, where $M$ is a Turing machine and $x$ a string}
  \Ensure{accept if $M(x)$ accepts, reject otherwise}
  \Function{$M_{\text{ACC}}$}{$\inner{M}, x$}
    \State Construct the code $\inner{M'}$ of the following machine $M'$:
    \Statex \hspace{2em} \textbf{function} $M'(w)$: \Comment{$w$ is ignored}
    \Statex \hspace{4em} Run $M(x)$ \Comment{$\inner{M}$ and $x$ are hard-coded into $M'$}
    \Statex \hspace{4em} \textbf{if} $M(x)$ accepts \textbf{then return} accept \textbf{else return} reject
    \State \Return $O(\inner{M'}, \inner{M'})$
  \EndFunction
\end{algorithmic}

\textbf{$M_{\text{ACC}}$ always halts.}
Writing down $\inner{M'}$ is string manipulation (a fixed template with $\inner{M}$ and $x$ pasted in), and $M_{\text{ACC}}$ never runs $M$ or $M'$.
It then makes one call to $O$, which halts because $O$ decides $L$, and returns that answer.

\textbf{Key property of $M'$.}
As in part (a), $M'(w)$ ignores $w$ and does what $M(x)$ does: it accepts if $M(x)$ accepts, rejects if $M(x)$ rejects, and loops if $M(x)$ loops.
So $\varepsilon \in L(M')$ if and only if $M(x)$ accepts, and since $L(M') \cup L(M') = L(M')$,
\[
  \varepsilon \in L(M') \cup L(M') \iff M(x) \text{ accepts}.
\]

\textbf{Correctness.}
By the definition of $\Lacc$, the key property, the definition of $L$, the fact that $O$ decides $L$, and the code of $M_{\text{ACC}}$,
\begin{align*}
  (\inner{M}, x) \in \Lacc &\Longrightarrow M(x) \text{ accepts} \\
  &\Longrightarrow \varepsilon \in L(M') = L(M') \cup L(M') \\
  &\Longrightarrow (\inner{M'}, \inner{M'}) \in L \\
  &\Longrightarrow O(\inner{M'}, \inner{M'}) \text{ accepts} \\
  &\Longrightarrow M_{\text{ACC}} \text{ accepts}, \\[0.5em]
  (\inner{M}, x) \notin \Lacc &\Longrightarrow M(x) \text{ rejects or loops} \\
  &\Longrightarrow M'(\varepsilon) \text{ does not accept, so } \varepsilon \notin L(M') \cup L(M') \\
  &\Longrightarrow (\inner{M'}, \inner{M'}) \notin L \\
  &\Longrightarrow O(\inner{M'}, \inner{M'}) \text{ rejects} \\
  &\Longrightarrow M_{\text{ACC}} \text{ rejects}.
\end{align*}
So $M_{\text{ACC}}$ decides $\Lacc$ with an oracle for $L$, i.e., $\Lacc \lT L$.
A decider for $L$ would therefore give a decider for $\Lacc$, which does not exist, so $L$ is undecidable.

(One could also reduce from part (a): $\Leacc \lT L$ via the one-line machine that returns $O(\inner{M}, \inner{M})$.
I reduced from $\Lacc$ directly so that the starting point is a language proved undecidable in lecture.)

\qpage{Q4.1 HALTs (a)}

$\Lbh$ is \textbf{decidable}.

The point is that both parts of the condition are bounded by numbers we can read off the input.
Let $\Sigma$ be the input alphabet of $M$ (it is listed in $\inner{M}$).
There are only finitely many strings over $\Sigma$ of length at most $k$, namely $\sum_{\ell=0}^{k} \abs{\Sigma}^\ell$ of them, and for each one we only need to watch $M$ for $376$ steps.
So we can just try them all.

\begin{algorithmic}[1]
  \Require{$(\inner{M}, k)$, where $M$ is a Turing machine and $k \ge 0$ an integer}
  \Ensure{accept if some input $x$ with $\abs{x} \le k$ has $M(x)$ halting within $376$ steps, reject otherwise}
  \Function{$D$}{$\inner{M}, k$}
    \For{$\ell = 0, 1, \ldots, k$}
      \For{each string $x \in \Sigma^\ell$ (in alphabetical order)}
        \State Simulate $M$ on $x$ for at most $376$ steps
        \If{$M$ entered its accept or reject state during these steps} \Return accept
        \EndIf
      \EndFor
    \EndFor
    \State \Return reject
  \EndFunction
\end{algorithmic}

\textbf{$D$ always halts.}
The two loops go over finitely many strings: $\abs{\Sigma}^\ell$ strings for each $\ell \le k$, and $k$ is a fixed number given in the input.
Each simulation is cut off after at most $376$ steps of $M$, and simulating one step of $M$ (look up the transition for the current state and symbol, update the tape and head) takes finite time.
So every iteration takes finitely many steps, there are finitely many iterations, and $D$ returns.
$D$ never runs $M$ without a step bound.

\textbf{Correctness.}
The simulation is exact, so for each $x$, $M$ halts within the first $376$ simulated steps if and only if $M(x)$ really halts in at most $376$ steps.
\begin{itemize}
  \item If $(\inner{M}, k) \in \Lbh$, there is some $x^*$ with $\abs{x^*} \le k$ such that $M(x^*)$ halts in at most $376$ steps.
  The loops eventually reach $x^*$ (unless $D$ has already accepted on an earlier string, which is also fine), the simulation of $M(x^*)$ sees $M$ halt, and $D$ accepts.
  \item If $(\inner{M}, k) \notin \Lbh$, then for every $x$ with $\abs{x} \le k$, $M(x)$ does not halt within $376$ steps.
  So no simulation sees a halting state, the \textbf{if} never fires, the loops run to completion, and $D$ rejects.
\end{itemize}
So $D$ decides $\Lbh$, and $\Lbh$ is decidable.

(In $376$ steps $M$ can only read the first $376$ cells of its tape, so it would be enough to try strings of length at most $\min(k, 376)$, but the search above is already finite, so I did not bother.)

\qpage{Q4.2 HALTs (b)}

$\Llh$ is \textbf{undecidable}.
I show $\Lh \lT \Llh$; since $\Lh$ is undecidable (lecture), so is $\Llh$.

Unlike part (a), the step budget $\abs{w}$ is not fixed: it grows with the input.
So a machine that takes the same number $T$ of steps on every input halts ``in time'' on every input of length at least $T$.

Let $O$ be an oracle that decides $\Llh$: $O$ halts on every input, and $O(\inner{N})$ accepts if and only if there is some string $w$ such that $N(w)$ halts in at most $\abs{w}$ steps.

\begin{algorithmic}[1]
  \Require{$(\inner{M}, x)$, where $M$ is a Turing machine and $x$ a string}
  \Ensure{accept if $M(x)$ halts, reject otherwise}
  \Function{$M_{\text{HALT}}$}{$\inner{M}, x$}
    \State Construct the code $\inner{M'}$ of the following machine $M'$:
    \Statex \hspace{2em} \textbf{function} $M'(w)$: \Comment{$w$ is never read}
    \Statex \hspace{4em} Run $M(x)$ \Comment{$\inner{M}$ and $x$ are hard-coded into $M'$}
    \Statex \hspace{4em} \textbf{return} accept
    \State \Return $O(\inner{M'})$
  \EndFunction
\end{algorithmic}

\textbf{$M_{\text{HALT}}$ always halts.}
Writing down $\inner{M'}$ is string manipulation, and $M_{\text{HALT}}$ never runs $M$ or $M'$.
It then makes one call to $O$, which halts because $O$ decides $\Llh$, and returns that answer.

\textbf{Key property of $M'$.}
$M'$ never looks at its input, so its computation is literally the same on every input $w$.
(It must not spend $\abs{w}$ steps erasing $w$ first, and it does not need to: $M'$ writes $x$ over the first $\abs{x}$ cells using a marked copy of $M$'s tape alphabet, and during the simulation treats every unmarked cell, leftover symbol of $w$ or blank, as a blank.)
So if $M(x)$ halts, there is a number $T \ge 1$, not depending on $w$, such that $M'(w)$ halts after exactly $T$ steps for every $w$; and if $M(x)$ loops, then $M'(w)$ loops for every $w$.

\textbf{Correctness.}
Let $a$ be any symbol of the input alphabet.
By the definition of $\Lh$, the key property, the definition of $\Llh$, the fact that $O$ decides $\Llh$, and the code of $M_{\text{HALT}}$,
\begin{align*}
  (\inner{M}, x) \in \Lh &\Longrightarrow M(x) \text{ halts} \\
  &\Longrightarrow M' \text{ halts on every } w \text{ in exactly } T \text{ steps, so on } w = a^T \text{ in } T = \abs{w} \text{ steps} \\
  &\Longrightarrow \inner{M'} \in \Llh \\
  &\Longrightarrow O(\inner{M'}) \text{ accepts} \\
  &\Longrightarrow M_{\text{HALT}} \text{ accepts}, \\[0.5em]
  (\inner{M}, x) \notin \Lh &\Longrightarrow M(x) \text{ loops} \\
  &\Longrightarrow M' \text{ loops on every } w\text{, so no } w \text{ has } M' \text{ halting within } \abs{w} \text{ steps} \\
  &\Longrightarrow \inner{M'} \notin \Llh \\
  &\Longrightarrow O(\inner{M'}) \text{ rejects} \\
  &\Longrightarrow M_{\text{HALT}} \text{ rejects}.
\end{align*}
So $M_{\text{HALT}}$ decides $\Lh$ using an oracle for $\Llh$, i.e., $\Lh \lT \Llh$.
A decider for $\Llh$ would give a decider for $\Lh$, which does not exist, so $\Llh$ is undecidable.

\qpage{Q5.1 Recognizing Unrecognizability (a)}

$\Lloops$ is \textbf{unrecognizable}.

The reason is that $\Lloops$ is essentially the complement of $\Lh$, and $\Lh$ is recognizable but not decidable.
We saw in lecture that if both a language and its complement are recognizable, then the language is decidable.
I include the argument below, since it is the whole proof.

\textbf{Setup.}
Let $V$ be the set of strings that are valid encodings of a pair $(\inner{M}, x)$ with $M$ a Turing machine and $x$ a string.
$V$ is decidable by a syntactic check, as in Q2(c).
For every pair in $V$, $M(x)$ either halts or loops, and not both, so
\[
  \Lh \cup \Lloops = V \qquad\text{and}\qquad \Lh \cap \Lloops = \emptyset .
\]
Also, $\Lh$ is recognizable (lecture): the universal machine $U(\inner{M}, x)$ runs $M$ on $x$ and accepts as soon as it halts, so $U$ accepts exactly the pairs in $\Lh$.

\textbf{Suppose for contradiction that $\Lloops$ is recognizable}, and let $R$ be a machine that accepts exactly the strings in $\Lloops$.
Consider the following machine.

\begin{algorithmic}[1]
  \Require{A string $w$}
  \Ensure{accept if $w \in \Lh$, reject otherwise}
  \Function{$D$}{$w$}
    \If{$w \notin V$} \Return reject \Comment{not a valid pair; decidable check}
    \EndIf
    \Loop
      \State Run $U(w)$ for one more step; \textbf{if} $U$ has accepted \textbf{then return} accept
      \State Run $R(w)$ for one more step; \textbf{if} $R$ has accepted \textbf{then return} reject
    \EndLoop
  \EndFunction
\end{algorithmic}
(Both simulations keep their own tape and state between iterations, so this alternates steps of $U(w)$ and $R(w)$; if one of them halts without accepting, $D$ just keeps stepping the other one.)

\textbf{$D$ always halts.}
If $w \notin V$, $D$ rejects right away.
If $w \in V$, then $w \in \Lh$ or $w \in \Lloops$.
In the first case $U(w)$ accepts after some finite number $t$ of steps, so $D$ returns by the end of iteration $t$.
In the second case $R(w)$ accepts after some finite number of steps, so again $D$ returns.
Either way the loop stops.

\textbf{Correctness.}
\begin{itemize}
  \item If $w \in \Lh$, then $w \in V$, $U(w)$ accepts, and $R(w)$ never accepts (because $w \notin \Lloops$ and $R$ recognizes $\Lloops$).
  So the only way out of the loop is ``return accept,'' and $D$ accepts.
  \item If $w \notin \Lh$, then either $w \notin V$ and $D$ rejects on line 1, or $w \in V$ with $M(x)$ looping, i.e., $w \in \Lloops$.
  Then $R(w)$ accepts and $U(w)$ never accepts, so $D$ rejects.
\end{itemize}
So $D$ decides $\Lh$.
But $\Lh$ is undecidable (lecture).
This contradiction shows that no such $R$ exists: $\Lloops$ is unrecognizable.

\qpage{Q5.2 Recognizing Unrecognizability (b)}

$\LSp$ is \textbf{recognizable} (for every fixed finite set $S$).

Write $S = \set{s_1, s_2, \ldots, s_n}$; since $S$ is fixed and finite, this list can be hard-coded into the recognizer.
Note $\inner{M} \in \LSp$ if and only if $M$ accepts every one of $s_1, \ldots, s_n$.
The recognizer just runs $M$ on each $s_i$ in turn.
It may get stuck forever on some $s_i$, but that only happens when $M$ does not accept $s_i$, so it is allowed.

\begin{algorithmic}[1]
  \Require{A string $w$}
  \Ensure{accept if $w = \inner{M}$ with $L(M) \supseteq S$; otherwise reject or loop}
  \Function{$R$}{$w$}
    \If{$w$ is not a valid encoding $\inner{M}$ of a Turing machine} \Return reject
    \EndIf
    \For{$i = 1, 2, \ldots, n$}
      \State Run $M(s_i)$ \Comment{this may loop, which is allowed for a recognizer}
      \If{$M(s_i)$ rejects} \Return reject
      \EndIf
    \EndFor
    \State \Return accept \Comment{every $s_i$ was accepted}
  \EndFunction
\end{algorithmic}

A recognizer does not have to halt on every input, so what has to be shown is that $R$ accepts $w$ if and only if $w \in \LSp$.

\textbf{If $w \in \LSp$, then $R$ accepts $w$.}
Then $w = \inner{M}$ for a Turing machine $M$ with $L(M) \supseteq S$, so the validity check passes, and $M$ accepts every $s_i$.
So in iteration $i$ the simulation of $M(s_i)$ halts, in the accept state, and the \textbf{if} does not fire.
After all $n$ iterations, $R$ reaches the last line and accepts.
(If $S = \emptyset$, the loop is empty and $R$ accepts every valid $\inner{M}$, which is right, since every $L(M) \supseteq \emptyset$.)

\textbf{If $w \notin \LSp$, then $R$ does not accept $w$.}
If $w$ is not a valid encoding, $R$ rejects on line 1.
Otherwise $w = \inner{M}$ and $L(M) \not\supseteq S$, so there is some $s_j \in S$ that $M$ does not accept.
Take the smallest such $j$.
For $i < j$, $M(s_i)$ accepts, so iterations $1, \ldots, j-1$ finish without returning.
In iteration $j$, $M(s_j)$ either rejects, in which case $R$ rejects, or loops, in which case $R$ loops inside the simulation and never reaches the accept line.
Either way $R$ does not accept.

So $R$ accepts exactly the strings in $\LSp$, i.e., $R$ recognizes $\LSp$.

(The order in which the $s_i$ are tried does not matter, and running the $n$ simulations in parallel would also work.
Running them one after another is simpler and is enough, because the only inputs where $R$ gets stuck are inputs it is supposed to not accept anyway.)

\qpage{Q6 Lagging Barber}

Recall $\Lb = \set{\inner{M} : M(\inner{M}) \text{ does not accept}}$, which is undecidable (lecture).
I show $\Lb \lT \Llag$, so $\Llag$ is undecidable as well.

Let $O$ be an oracle that decides $\Llag$: $O$ halts on every input, and $O(\inner{M_1}, \inner{M_2})$ accepts if and only if for every $x \in \Sigma^*$, $M_2(x)$ accepts implies $M_1(x)$ loops.

\textbf{Idea.}
Choose $M_1$ to be a machine that \emph{never} loops.
Then ``$M_2(x)$ accepts $\Rightarrow$ $M_1(x)$ loops'' can only hold when $M_2(x)$ does not accept, so $M_1$ lags behind $M_2$ exactly when $M_2$ accepts nothing.
Now choose $M_2$ to accept everything if $M(\inner{M})$ accepts, and nothing otherwise; then the oracle answers exactly the Barber question.

\begin{algorithmic}[1]
  \Require{$\inner{M}$, where $M$ is a Turing machine}
  \Ensure{accept if $M(\inner{M})$ does not accept, reject otherwise}
  \Function{$\Mb$}{$\inner{M}$}
    \State Let $M_{\text{yes}}$ be the machine that accepts every input immediately (fixed code)
    \State Construct the code $\inner{M'}$ of the following machine $M'$:
    \Statex \hspace{2em} \textbf{function} $M'(w)$: \Comment{$w$ is ignored}
    \Statex \hspace{4em} Run $M(\inner{M})$ \Comment{$\inner{M}$ is hard-coded into $M'$}
    \Statex \hspace{4em} \textbf{if} $M(\inner{M})$ accepts \textbf{then return} accept \textbf{else return} reject
    \State \Return $O(\inner{M_{\text{yes}}}, \inner{M'})$ \Comment{$M_1 = M_{\text{yes}}$, $M_2 = M'$}
  \EndFunction
\end{algorithmic}

\textbf{$\Mb$ always halts.}
$\inner{M_{\text{yes}}}$ is a fixed string, and writing down $\inner{M'}$ is string manipulation with $\inner{M}$ pasted into a template.
$\Mb$ never runs $M$, $M'$, or $M_{\text{yes}}$.
It makes one call to $O$, which halts because $O$ decides $\Llag$, and returns that answer.

\textbf{Key properties.}
$M_{\text{yes}}(x)$ halts on every $x$, so $M_{\text{yes}}$ never loops.
$M'(w)$ ignores $w$ and does what $M(\inner{M})$ does, so either $M'$ accepts every string (if $M(\inner{M})$ accepts) or $M'$ accepts no string (if $M(\inner{M})$ rejects or loops).

\textbf{Correctness.}
By the definition of $\Lb$, the key properties, the definition of lagging, the fact that $O$ decides $\Llag$, and the code of $\Mb$,
\begin{align*}
  \inner{M} \in \Lb &\Longrightarrow M(\inner{M}) \text{ does not accept} \\
  &\Longrightarrow M'(x) \text{ does not accept, for every } x \in \Sigma^* \\
  &\Longrightarrow \text{``$M'(x)$ accepts $\Rightarrow$ $M_{\text{yes}}(x)$ loops'' holds vacuously for every } x \\
  &\Longrightarrow M_{\text{yes}} \text{ lags behind } M'\text{, i.e., } (\inner{M_{\text{yes}}}, \inner{M'}) \in \Llag \\
  &\Longrightarrow O(\inner{M_{\text{yes}}}, \inner{M'}) \text{ accepts} \\
  &\Longrightarrow \Mb \text{ accepts}, \\[0.5em]
  \inner{M} \notin \Lb &\Longrightarrow M(\inner{M}) \text{ accepts} \\
  &\Longrightarrow M'(\varepsilon) \text{ accepts, but } M_{\text{yes}}(\varepsilon) \text{ halts rather than loops} \\
  &\Longrightarrow \text{the lag condition fails at } x = \varepsilon\text{, i.e., } (\inner{M_{\text{yes}}}, \inner{M'}) \notin \Llag \\
  &\Longrightarrow O(\inner{M_{\text{yes}}}, \inner{M'}) \text{ rejects} \\
  &\Longrightarrow \Mb \text{ rejects}.
\end{align*}
So $\Mb$ decides $\Lb$ using an oracle for $\Llag$, i.e., $\Lb \lT \Llag$.
If $\Llag$ were decidable, replacing $O$ by its decider would give a decider for $\Lb$, which does not exist.
Therefore $\Llag$ is undecidable.

\qpage{Q7.1 Infinite machines, infinite inputs! (a)}

The flaw is that the test for $M_1$ never finishes if $M_1$ does not halt on any input, so the procedure never gets to $M_2$.

Dovetailing over the inputs $s_1, s_2, s_3, \ldots$ of $M_1$ does fix one problem: no single looping computation can block the others, since every $M_1(s_j)$ keeps getting steps.
And if $M_1$ halts on some $s_j$, the dovetailing will eventually see that, because after enough rounds $M_1(s_j)$ has been given as many steps as it needs.
But there are infinitely many inputs.
If $M_1$ loops on every one of them, then no round ever produces a halt, and the dovetailing just keeps starting new computations $M_1(s_4), M_1(s_5), \ldots$ forever.
There is no point at which Professor $\Upsilon$ can declare ``$M_1$ is done, move on to $M_2$,'' because ``$M_1$ halts on no input'' is not something that can be observed in finitely many steps.

So take a collection in which $M_1$ loops on everything (say $M_1$ immediately enters an infinite loop), while $M_2$ halts on $s_1$ in one step.
The correct answer is ``yes, some machine halts on some input,'' but Professor $\Upsilon$'s procedure spends forever on $M_1$ and never runs a single step of $M_2$.
Dovetailing has to cover \emph{both} infinite dimensions, machines and inputs, at the same time; dovetailing over one of them and then iterating over the other does not work.

(There is a smaller issue in the phrasing as well: Professor $\Upsilon$ talks about testing whether $M_1$ ``halts on all its possible inputs,'' but the goal only needs some machine to halt on some input, which is what the dovetailing can detect.
Halting on \emph{all} inputs is not something any finite amount of simulation can confirm either.)

\qpage{Q7.2 Infinite machines, infinite inputs! (b)}

The flaw is the same one, with the roles of machines and inputs swapped: the test for $s_1$ never finishes if no machine halts on $s_1$, so the procedure never gets to $s_2$.

Dovetailing over the machines $M_1, M_2, M_3, \ldots$ on the fixed input $s_1$ makes sure that every $M_i(s_1)$ keeps getting steps, so if some $M_i$ halts on $s_1$, that will eventually be seen.
But there are infinitely many machines.
If every $M_i$ loops on $s_1$, then no round ever produces a halt, and Dr.\ Epsiwrong's procedure keeps adding $M_4(s_1), M_5(s_1), \ldots$ to the pile forever.
The event ``no machine halts on $s_1$'' cannot be observed after finitely many steps, so there is never a moment to say ``done with $s_1$, move on to $s_2$.''

Concretely, take a collection in which every $M_i$ loops on $s_1$ (say each $M_i$ loops whenever its input is $s_1$), while $M_1$ halts on $s_2$ in one step.
The correct answer is ``yes,'' but the procedure spends forever on $s_1$ and never runs a single step on $s_2$.

So both colleagues made the same mistake.
Each of them dovetails along one infinite list and then tries to iterate over the other infinite list, and the iteration never advances past its first element.
The fix, in part (c), is to dovetail over machines and inputs simultaneously, so that every pair $(M_i, s_j)$ gets started after finitely many rounds and then gets a step in every round after that.

\qpage{Q7.3 Infinite machines, infinite inputs! (c)}

\textbf{Setup.}
Let $s_1, s_2, s_3, \ldots$ be the standard enumeration of $\Sigma^*$ (by length, and alphabetically within each length); given $j$, the string $s_j$ can be computed.
I assume the archive gives us access to the machines in order, i.e., given $i$ we can look up the code $\inner{M_i}$ (for example, the encodings are stored one after another on a read-only tape).
Without some such access there is nothing to simulate.
The machine below ignores its own input.

\begin{algorithmic}[1]
  \Require{Any string (ignored)}
  \Ensure{accept if some $M_i$ halts on some $s_j$; loop otherwise}
  \Function{$A$}{}
    \For{$n = 1, 2, 3, \ldots$} \Comment{round $n$}
      \For{each pair $(i, j)$ with $1 \le i \le n$ and $1 \le j \le n$}
        \If{the pair $(i,j)$ has not been started yet} look up $\inner{M_i}$ and $s_j$ and set up the initial configuration of $M_i$ on $s_j$
        \EndIf
        \State Simulate one more step of $M_i(s_j)$, continuing from its saved configuration
        \If{$M_i(s_j)$ has halted (accept or reject)} \Return accept
        \EndIf
      \EndFor
    \EndFor
  \EndFunction
\end{algorithmic}
At any time the machine stores the configurations of the finitely many pairs started so far (at most $n^2$ of them during round $n$), so each round is a finite amount of work.

\textbf{When a pair gets its steps.}
The pair $(i, j)$ first appears in round $\max(i, j)$ and gets exactly one step in that round and in every later round.
So by the end of round $\max(i, j) + t - 1$, the simulation of $M_i(s_j)$ has been run for exactly $t$ steps (unless $A$ accepted earlier).

\textbf{Each round finishes.}
Round $n$ handles $n^2$ pairs, and for each one it does one lookup (only the first time), one simulated step, and one check.
All of these take finite time, so round $n$ ends after finitely many steps and round $n+1$ begins, unless $A$ returns during round $n$.

\textbf{Correctness.}
\begin{itemize}
  \item Suppose some machine in the list halts on some input: $M_i(s_j)$ halts, say after $t$ steps.
  Since every round finishes, $A$ reaches round $\max(i,j) + t - 1$ unless it has accepted before that.
  By then $M_i(s_j)$ has been simulated for $t$ steps, so the simulation has reached a halting state, and the check after that step makes $A$ accept.
  So $A$ accepts.
  \item Suppose no machine in the list halts on any input.
  Then no simulated computation $M_i(s_j)$ ever reaches a halting state, so the check never fires and ``return accept'' is never executed.
  Every round finishes and the next one starts, so $A$ runs through rounds $1, 2, 3, \ldots$ forever.
  So $A$ loops.
\end{itemize}
Hence $A$ accepts if and only if some $M_i$ halts on some $s_j$, and loops otherwise, as required.

(Note that $A$ is a recognizer, not a decider: it never answers ``no.''
That is all the problem asks for.)

\qpage{Q8.1 (Optional extra credit) (a)}

\textbf{True.} For every undecidable $L_1$ and finite $L_2$, $L_1 \cup L_2$ is undecidable.

The idea: adding finitely many strings cannot make an undecidable language decidable, because a decider for the union, together with a finite table of answers for those finitely many strings, would decide $L_1$.
I write this as $L_1 \lT L_1 \cup L_2$.

Let $L_2 = \set{y_1, \ldots, y_m}$, and let $T = L_1 \cap L_2$.
$T$ is a subset of a finite set, so $T$ is finite, and a finite language is decided by a machine that hard-codes its members.
We do not need to know which $y_i$ are in $T$; it is enough that the finite list exists, so the machine $D_T$ deciding $T$ exists (this is the same kind of argument as for $\Lsmh$ in HW6).
Let $O$ be an oracle that decides $L_1 \cup L_2$.

\begin{algorithmic}[1]
  \Require{A string $x$}
  \Ensure{accept if $x \in L_1$, reject otherwise}
  \Function{$M_1$}{$x$}
    \If{$x \in \set{y_1, \ldots, y_m}$} \Return $D_T(x)$ \Comment{finite table lookup}
    \Else{} \Return $O(x)$
    \EndIf
  \EndFunction
\end{algorithmic}

\textbf{$M_1$ always halts.}
Comparing $x$ against the $m$ fixed strings is finite work, $D_T$ halts on every input because it is a decider for a finite language, and $O$ halts on every input because it decides $L_1 \cup L_2$.

\textbf{Correctness.}
\begin{itemize}
  \item If $x \in L_2$: then $x \in L_1 \iff x \in L_1 \cap L_2 = T$, and $D_T$ decides $T$, so $M_1(x)$ accepts if $x \in L_1$ and rejects otherwise.
  \item If $x \notin L_2$: then $x \in L_1 \iff x \in L_1 \cup L_2$ (a string outside $L_2$ is in the union exactly when it is in $L_1$), and $O$ decides $L_1 \cup L_2$, so again $M_1(x)$ accepts if $x \in L_1$ and rejects otherwise.
\end{itemize}
So $M_1$ decides $L_1$ with an oracle for $L_1 \cup L_2$, i.e., $L_1 \lT L_1 \cup L_2$.
If $L_1 \cup L_2$ were decidable, replacing $O$ by its decider would give a decider for $L_1$, contradicting the assumption that $L_1$ is undecidable.
Hence $L_1 \cup L_2$ is undecidable.

\qpage{Q8.2 (Optional extra credit) (b)}

\textbf{False.} An undecidable language can be made decidable by taking a union with a decidable one.

Counterexample: $L_1 = \Lh$ and $L_2 = \Sigma^*$.
\begin{itemize}
  \item $L_1 = \Lh$ is undecidable (lecture).
  \item $L_2 = \Sigma^*$ is decidable: the machine that accepts every input immediately halts on every input and accepts exactly the strings in $\Sigma^*$.
  \item $L_1 \cup L_2 = \Lh \cup \Sigma^* = \Sigma^*$, which is decidable by the same machine.
\end{itemize}
So here $L_1$ is undecidable, $L_2$ is decidable, and $L_1 \cup L_2$ is decidable.
This disproves the statement.

This does not contradict part (a).
The argument there used that $L_2$ is finite, so that the strings where the union does not determine $L_1$ could be handled by a finite table.
When $L_2$ is infinite, the union may hide infinitely much information about $L_1$, as it does here: $\Sigma^*$ swallows $\Lh$ completely.

\end{document}
